\section{A normalized factorial on a complex strip}\label{ca:factorial-section}

We first fix the interpolation used in every subsequent contour. Put
$Q=e^{2\pi\ii/N}$ and
\[
 S_r=\prod_{j=1}^r2\sin(\pi j/N),\qquad S_0=1.
\]
The function below is a normalized solution of a sine difference equation. Only the ordinary Gamma functional equation, its Stirling expansion with remainder, and the convergent sine product are needed; the standard Gamma formulas are recorded in \cite[\S\S5.5(i), 5.11(i)--(ii)]{DLMF}.

\begin{lemma}[Canonical factorial interpolation]\label{ca:factorial}
For $N\ge2$ and $-1<\Re z<N$ the product
\begin{equation}\label{ca:G}
 G_N(z)=\left(\frac{2\pi}{N}\right)^z\Gamma(z+1)
 \prod_{k=1}^{\infty}
 \frac{\Gamma(Nk+z+1)}{\Gamma(Nk-z)(Nk)^{2z+1}}
\end{equation}
is holomorphic and nonzero. Positive-base powers use their real logarithms. With $d=N-1$ and
\[
 \widehat F_N(z)=e^{\pi\ii z(z+1)/(2N)-\pi\ii z/2}G_N(z)
\]
one has the exact identities
\begin{align}
 G_N(z+1)&=2\sin\!\frac{\pi(z+1)}N\,G_N(z),\notag\\
 \widehat F_N(z+1)&=(1-Q^{z+1})\widehat F_N(z),\label{ca:shift}\\
 G_N(0)&=1,\qquad G_N(r)=S_r,\qquad
 \widehat F_N(r)=(Q;Q)_r\quad(0\le r<N),\notag\\
 G_N(z)G_N(d-z)&=N,\notag\\
 \widehat F_N(z)\widehat F_N(d-z)&=
 N e^{\pi\ii z(z-d)/N}.\label{ca:reflection}
\end{align}
Identities are asserted where the factors are regular, and elsewhere by meromorphic continuation. If
\[
 H_N(s)=\prod_{j=1}^{d}
       \frac{\sin(\pi(j-s)/N)}{\sin(\pi j/N)}
       =\frac{\sin(\pi s)}{N\sin(\pi s/N)},
\]
then on $|s|<1$,
\begin{equation}\label{ca:small-reflection}
 G_N(s)G_N(-s)=H_N(s)^{-1}.
\end{equation}

Let $J$ be the analytic primitive of $\Log(2\sin\pi x)$ on
$0<\Re x<1$, fixed by its integrable limit $J(0)=0$ along the real interval. On every compact subset of this strip,
\begin{equation}\label{ca:Stirling}
 \log G_N(Nx)
 =NJ(x)+\tfrac12\Log(2N\sin\pi x)
       +\frac{\pi\cot\pi x}{12N}+O(N^{-3}).
\end{equation}
The branches are real on $0<x<1$. More generally, for each fixed $L\ge1$ the correction is
\begin{align}
 &\sum_{r=1}^{L}\frac{c_{2r-1}(x)}{N^{2r-1}}+O(N^{-2L-1}),\notag\\
 &c_{2r-1}(x)=\frac{B_{2r}}{2r(2r-1)}
 \left\{x^{1-2r}+
 \sum_{k\ge1}\bigl[(k+x)^{1-2r}-(k-x)^{1-2r}\bigr]\right\}.
 \label{ca:all-Stirling}
\end{align}
Here $B_j$ are defined by $t/(e^t-1)=\sum B_jt^j/j!$.
All fixed derivatives have the same uniform remainder orders on smaller compact sets. Fixed bounded complex offsets are included by substituting $x+c/N$; in particular
\begin{equation}\label{ca:offset}
 G_N(Nx+c)=
 \sqrt{2N\sin\pi x}\,e^{NJ(x)}(2\sin\pi x)^c
 \bigl(1+O(N^{-1})\bigr).
\end{equation}
\end{lemma}

\begin{proof}
The logarithm of the $k$th factor in \eqref{ca:G} is
\[
 -\frac{z(z+1)(2z+1)}{6N^2k^2}+O(k^{-3})
\]
locally uniformly in $z$. The product and all its local logarithmic derivatives therefore converge. Its factors are nonzero in the stated strip and its value at zero is one. Gamma's functional equation gives
\[
 \frac{G_N(z+1)}{G_N(z)}
 =\frac{2\pi(z+1)}N
  \prod_{k\ge1}\left(1-\frac{(z+1)^2}{N^2k^2}\right)
 =2\sin\frac{\pi(z+1)}N.
\]
Multiplying by the displayed exponential gives the second shift. This also fixes the lattice values, including their phases.

To remove any periodic ambiguity in reflection, differentiate
$\log[G_N(z)G_N(d-z)]$. In the product truncated at $k=K$ the resulting sum telescopes to
\[
 \psi(NK+z+1)-\psi(N(K+1)-z),
\]
where $\psi=\Gamma'/\Gamma$. It tends locally uniformly to zero. The product is constant, and at $z=0$ it equals $S_d=N$. The latter equality follows by taking the absolute value of
$\prod_{j=1}^{N-1}(1-e^{2\pi\ii j/N})=N$.
This proves both formulas in \eqref{ca:reflection}, with the marked phase obtained by adding the two quadratic exponents. Iterating the shift from $-s$ to $d-s$ gives
$G_N(d-s)=N H_N(s)G_N(-s)$, which proves \eqref{ca:small-reflection}. The finite sine product gives the displayed quotient for $H_N$, with its removable value one at zero.

We give the remainder argument for \eqref{ca:Stirling}, because contours require a complex estimate rather than values on an integer mesh. Stirling's expansion on a closed right-half-plane sector has remainder $O(|w|^{-2L-1})$ after its term in $w^{1-2L}$. One obtains this bound from the Euler summation formula for $\log\Gamma$: the remaining periodic Bernoulli polynomial is bounded, and its integral is bounded by a constant times
$\int_0^\infty|w+t|^{-2L-2}\dd t$.
For $x$ in a fixed complex compact set as above, $Nx$ and
$N(k\pm x)$ lie in a common such sector.

For $L=1$ the logarithm of the $k$th factor in \eqref{ca:G}, at $z=Nx$, becomes
\[
 N\phi_k(x)+\tfrac12\Log(1-x^2/k^2)
 +\frac1{12N}\left(\frac1{k+x}-\frac1{k-x}\right)
 +O(N^{-3}k^{-3}),
\]
where
\[
 \phi_k(x)=(k+x)\Log(1+x/k)
 -(k-x)\Log(1-x/k)-2x.
\]
The remainder is summable. Moreover $\phi_k(0)=0$ and
$\phi_k'(x)=\Log(1-x^2/k^2)$. The sine product, including the
$\Gamma(Nx+1)$ factor, therefore sums the leading terms to $NJ(x)$.
It sums the half logarithms to $\frac12\Log(2N\sin\pi x)$.
The identity
\[
 \frac1x+\sum_{k\ge1}
   \left(\frac1{k+x}-\frac1{k-x}\right)=\pi\cot\pi x
\]
gives the first correction. Keeping further Stirling terms gives
\eqref{ca:all-Stirling}; their paired series converge locally uniformly, and the remainder sums absolutely as
$O(N^{-2L-1}\sum k^{-2L-1})$. The normalization of $J$, including its additive constant, was fixed by integration from zero, not by differentiation alone. Cauchy's formula on a larger compact set proves the derivative assertions. Taylor expansion at $x+c/N$ proves \eqref{ca:offset}.
\end{proof}

The reflection $J(1-x)=-J(x)$ follows by differentiating and using
$\int_0^1\log(2\sin\pi t)\dd t=0$. In particular $J(1/2)=0$. We will use
\[
 J(1/2+t)=-J(1/2-t)
\]
to normalize a sheared action at the actual terminal index.
